\documentclass[10pt,a4paper]{amsart}
\usepackage[margin=19mm]{geometry}
\usepackage{amsmath,amssymb,mathtools}
\usepackage[scr=rsfs,cal=euler]{mathalfa}
\usepackage{xfrac}
\usepackage[unicode,hidelinks,bookmarks=false]{hyperref}
\newcommand{\ie}{i.e.\ }
\newcommand{\cf}{cf.\ }
\newcommand{\resp}{respectively\ }
\newcommand{\Subsection}{\subsection}
\providecommand{\nobreakdash}{\nobreak\hbox{-}}
\setlength{\emergencystretch}{2em}
\multlinegap0pt
\usepackage{amssymb}
\usepackage{mathtools}
\usepackage{xcolor}
\usepackage{algorithm}\usepackage{algpseudocode}
\newtheorem{assumption}{Assumption}
\newtheorem{theorem}{Theorem}
\usepackage{cleveref}
\crefname{assumption}{Assumption}{Assumptions}
\crefname{theorem}{Theorem}{Theorems}
\newcommand{\Prob}{\mathbb{P}}
\newcommand{\R}{\mathbb{R}}
\newcommand{\calU}{\mathcal{U}}
\newcommand{\vstar}{v^{\star}}
\title{Controlling the Swarm: Sparse Actuation and Collision Avoidance under Stochastic Delay}
\author{Jiguang Yu}
\date{Source: arXiv:2605.00395v1 (CC BY 4.0)}
\begin{document}
\maketitle

\noindent\emph{Source and licence.} Controlling the Swarm: Sparse Actuation and Collision Avoidance under Stochastic Delay. Version arXiv:2605.00395v1 (CC BY 4.0), distributed by arXiv under the Creative Commons Attribution 4.0 International licence, http://creativecommons.org/licenses/by/4.0/, as recorded in the arXiv licence metadata for this exact version and shown on the abstract page https://arxiv.org/abs/2605.00395v1. Full licence notice: \texttt{licenses/arxiv-2605.00395-CC-BY-4.0.txt}.\par\smallskip

\noindent\textit{Editorial scope.} Counted unit: the article's theorem thm:global (source0.tex lines 628--633). The article's complete original proof of it is delivered with it (source0.tex lines 634--639, one \texttt{proof} environment of 6 lines). No mathematics is written, repaired or re-derived: the statement and the proof are the article's own lines, in the article's own notation, and no retained line is substituted or re-flowed.  Retained uncounted as the source-local context the proof needs: lines 363--365; lines 367--376; lines 395--398; lines 557--566; lines 599--605; lines 612--616.  Nothing was omitted from any retained span, so the package carries no omission record.  No retained line contains a citation, so the wrapper carries no bibliography and no bibliography entry.  No retained line contains a figure environment, an image-inclusion command or drawing code, so no image is delivered and the package declares no asset dependency.  Editorial formatting only, in addition: an AMS \texttt{amsart} preamble that restates the article's own declarations and macros used by the retained lines, namely the environments \texttt{assumption}, \texttt{theorem} on separate counters, together with the standard packages those lines need.  The retained environments print in this document as Assumption 1, Theorem 1 in order of appearance, whereas the article prints the same items with the numbering of its own full document.  The complete native source of the article as distributed in the arXiv e-print for this exact version is bundled unchanged as \texttt{sources/199558/source0.tex}.
\begin{equation}\label{eq:pos}
  dx_i(t)=v_i(t)\,dt
\end{equation}

\begin{equation}\label{eq:vel}
\begin{aligned}
  dv_i(t) &= \Biggl[ \mathcal A_i^{\varepsilon,\delta}\bigl(t,\mathbf{x}(t),\mathbf{v}_t\bigr)
    - \frac{1}{N}\sum_{j\neq i}\nabla U\bigl(x_i(t)-x_j(t)\bigr) \\
  &\qquad - \nabla V_{\mathrm{obs}}\bigl(x_i(t)\bigr)
    - b_i\bigl(v_i(t)-\vstar\bigr) + B_i u_i(t) \Biggr]dt \\
  &\quad +\sigma_i\bigl(x_i(t),v_i(t),\mu_t^N\bigr)\,dW_i(t) \\
  &\quad +\sigma_i^{0}\bigl(x_i(t),v_i(t),\mu_t^N\bigr)\,dW^0(t),
\end{aligned}
\end{equation}

\begin{equation}\label{eq:compat}
  x_i^{\mathrm{in}}(t) = x_i^{\mathrm{in}}(0) - \int_t^0 v_i^{\mathrm{in}}(s)\,ds,
  \quad t\in[-\tau_{\max},0].
\end{equation}

\begin{assumption}[Regularized communication law]\label{ass:net}
The communication profile satisfies
$\phi\in C_b^1([0,\infty);[0,\infty))$. For each $i\neq j$, the regularized
coefficients $\chi_{ij}^{\varepsilon}:[0,\infty)\times(\R^d)^N\to[0,1]$ and
$r_{ij}^{\varepsilon}:(\R^d)^N\to[1,N-1]$ are bounded and locally Lipschitz
in $\mathbf{x}$, uniformly on compact time intervals. Consequently, the weights
$a_{ij}^{\varepsilon}(t,\mathbf{x}) = \chi_{ij}^{\varepsilon}(t,\mathbf{x})\,
\phi(r_{ij}^{\varepsilon}(\mathbf{x})/K)$ are also bounded and locally Lipschitz
in $\mathbf{x}$.
\end{assumption}

\begin{assumption}[Collision-free compatible history]\label{ass:hist}
The history paths
$x_i^{\mathrm{in}},v_i^{\mathrm{in}} \in C([-\tau_{\max},0];\R^d)$ satisfy
\eqref{eq:compat} and are collision-free:
$\inf_{s\in[-\tau_{\max},0]}\min_{i\neq j}
|x_i^{\mathrm{in}}(s)-x_j^{\mathrm{in}}(s)|>0$.
\end{assumption}

\begin{theorem}[Local strong well-posedness]\label{thm:local}
Suppose Assumptions~\ref{ass:net}--\ref{ass:hist} hold. For any $T>0$ and
$u\in\calU_T$, the system \eqref{eq:pos}--\eqref{eq:vel} admits a pathwise
unique strong solution on $[-\tau_{\max},\,T\wedge\tau_{\mathrm{coll}})$.
\end{theorem}

\begin{theorem}[Global strong well-posedness]\label{thm:global}
Assume the hypotheses of \Cref{thm:local} and that the repulsive interaction
ensures $\Prob(\tau_{\mathrm{coll}}=+\infty)=1$. Then, for every $T>0$, the
solution extends uniquely to a pathwise unique strong solution on
$[-\tau_{\max},T]$.
\end{theorem}

\begin{proof}[Proof sketch]
If $\Prob(\tau_{\mathrm{coll}}=+\infty)=1$, then $\tau_{\mathrm{coll}}>T$
almost surely for any $T>0$. Thus, the local solution never reaches the
singular set on $[0,T]$, allowing the unique extension throughout the
interval.
\end{proof}

\end{document}
